\documentclass[
  spanish,
  letterpaper, oneside
]{article}

\def\documenttitle {El interaprendizaje en ambientes virtuales}
\def\documentsubtitle {Lo que el interaprendizaje no es y cómo evitarlo}
\def\documentsubject {Actividad 3 - Interaprendizaje en ambientes virtuales}

\def\documentauthor {Martin Jonathan de la Cruz}
\def\coursename {Interaprendizaje en ambientes virtuales}
\def\coursecode {004}

\def\universityname {Universidad Abierta y a Distancia de México}
\def\universityfaculty {Licenciatura en Derecho}
\def\universitydepartment {Interaprendizaje en ambientes virtuales}
\def\universitydepartmentimage {departamentos/UnADM}
\def\universitydepartmentimagecfg {height=1.57cm}
\def\universitylocation {Roma Norte, Ciudad de México}

\def\authortable {
  \begin{tabular}{ll}
    \textbf{Alumno:}
    & \begin{tabular}[t]{l}
      Martin Jonathan de la Cruz \\
    \end{tabular} \\
    Matrícula: & ES2611202040 \\
    Figura docente: & Janeth Santana Ramírez \\
    Semestre/Bloque: & 2 / 1 \\
    & \\
    \multicolumn{2}{l}{Fecha de realización: \today} \\
    \multicolumn{2}{l}{\universitylocation}
  \end{tabular}
}

\input{template}
\usepackage{helvet}
\renewcommand{\familydefault}{\sfdefault}
\usepackage{array}
\usepackage{longtable}
\usepackage{pdflscape}
\usepackage{tikz}
\usetikzlibrary{arrows.meta,positioning,calc,matrix,fit,shapes.geometric,shadows.blur}
\setcitestyle{authoryear,open={(},close={)}}

\def\coverwatermarkenabled {true}
\def\coverwatermarkimage {img/departamentos/UnADM.pdf}
\def\coverwatermarkopacity {0.16}
\def\coverwatermarkwidth {0.72\paperwidth}
\def\coverwatermarkxshift {0cm}
\def\coverwatermarkyshift {-0.35cm}

\newcommand{\insertcoverwatermark}{%
  \ifthenelse{\equal{\coverwatermarkenabled}{true}}{%
    \AddToShipoutPictureBG*{%
      \begin{tikzpicture}[remember picture,overlay]
        \node[
          opacity=\coverwatermarkopacity,
          inner sep=0pt,
          xshift=\coverwatermarkxshift,
          yshift=\coverwatermarkyshift
        ] at (current page.center) {%
          \includegraphics[width=\coverwatermarkwidth]{\coverwatermarkimage}%
        };
      \end{tikzpicture}%
    }%
  }{}%
}

\definecolor{unadmGreenDark}{HTML}{174A3A}
\definecolor{unadmGreen}{HTML}{5F8F3A}
\definecolor{unadmGold}{HTML}{B88A2A}
\definecolor{unadmPaper}{HTML}{F6F7F2}
\definecolor{unadmInk}{HTML}{1F2A24}
\definecolor{unadmRowAlt}{HTML}{EEF2E6}

% -----------------------------------------------------------------------------
% AulaTeX:wiki-producto v1 -- bloque de contribución a la wiki colaborativa.
% La wiki es la herramienta con que se materializa la técnica SOCIOAPRENDIZAJE
% (#53 de las 100 Técnicas Didácticas UnADM): construcción coral del conocimiento.
% Caja institucional homóloga al forobox: banda verde, fondo claro, texto
% justificado. \wikimeta identifica la contribución publicada.
% -----------------------------------------------------------------------------
\usepackage[most]{tcolorbox}
\usepackage{attachfile}
\usepackage{enumitem}

% attachfile: el 'color' debe darse como TRIPLE RGB numérico (0..1); un nombre
% \definecolor[HTML] provoca 'Argument of \atfi@textcolor has an extra }'.
\attachfilesetup{color={0.373 0.561 0.227},print=false}

% Botón de copia: appearance={} oculta el PushPin y muestra el 2.o argumento como
% botón visible. Adjunta el .txt embebido con la contribución copiable.
\newcommand{\wikiCopyButton}[1]{%
  \textattachfile[appearance={},description={Contribucion a la wiki (texto copiable)},mimetype=text/plain]{#1}{%
    \colorbox{unadmGreen}{\small\textcolor{white}{\,Copiar contribución\,}}%
  }%
}

% Entorno visual de la wiki (mismo estilo que el forobox de la materia)
\newtcolorbox{wikibox}[1][]{%
  enhanced, breakable,
  colback=unadmPaper, colframe=unadmGreen,
  boxrule=0pt, leftrule=3pt, arc=1.2mm,
  left=4mm, right=4mm, top=3mm, bottom=3mm,
  fonttitle=\bfseries\color{white}, coltitle=white, colbacktitle=unadmGreenDark,
  attach boxed title to top left={xshift=4mm,yshift=-2mm},
  boxed title style={colframe=unadmGreenDark,arc=0.6mm},
  #1
}

% Ficha que identifica la contribución colaborativa (tema, participante, edición).
\newcommand{\wikimeta}[4]{%
  \begin{tcolorbox}[colback=unadmRowAlt,colframe=unadmRowAlt,boxrule=0pt,arc=1mm,
    left=2mm,right=2mm,top=1.5mm,bottom=1.5mm]
  {\small\begin{tabular}{@{}p{3.4cm}p{9.6cm}@{}}
    \textbf{Tema de la wiki} & #1 \\
    \textbf{Participante} & #2 \\
    \textbf{Fecha de edición} & #3 \\
    \textbf{Espacio / curso} & #4 \\
  \end{tabular}}
  \end{tcolorbox}\medskip}

\begin{document}

\insertcoverwatermark
\templatePortrait
\templatePagecfg
\onehalfspacing

\begin{abstractd}
  El interaprendizaje nombra la construcción conjunta del conocimiento entre pares:
  no se reduce a repartir tareas ni a acumular opiniones, sino que exige escucha
  activa, transformación de las ideas propias a partir de las ajenas y condiciones de
  participación que hagan de la diversidad un recurso. Sobre esa base conceptual
  —principios, condiciones y prácticas dialógicas frente a monológicas— este reporte
  contextualiza el trabajo colaborativo en ambientes virtuales y reúne los dos
  productos de la semana: una contribución en wiki que identifica una práctica
  contraria al interaprendizaje y el registro razonado del cuestionario
  ``Interaprendizaje: comprender, construir y distinguir''. El socioaprendizaje
  sostiene la construcción colaborativa, mientras el estudio de caso permite aplicar
  los conceptos a decisiones concretas de colaboración académica a distancia.
\end{abstractd}

\templateIndex
\templateFinalcfg

\section{Introducción}

Aprender en un ambiente virtual no equivale a estar conectados al mismo tiempo ni a
dividirnos el trabajo para terminar antes. El interaprendizaje parte de una premisa
distinta: el conocimiento se construye entre personas cuando cada quien aporta,
escucha y transforma sus propias ideas a partir de las de otras. Por eso, comprender
esta noción exige distinguir con claridad qué prácticas construyen conocimiento en
conjunto y cuáles solo lo simulan.

El propósito de este trabajo es doble. Primero, ordenar la base conceptual del
interaprendizaje —su noción, sus principios y condiciones, y la diferencia entre las
prácticas dialógicas y las monológicas— apoyándose en las lecturas de la semana.
Después, aplicar ese criterio en los dos productos previstos para la semana. La wiki
colaborativa funciona como espacio de construcción coral donde cada participante
añade sin borrar ni repetir lo aportado por otras personas. El cuestionario, por su
parte, verifica individualmente la comprensión y la transferencia de esos conceptos
a un caso de colaboración académica a distancia. En conjunto, ambos productos
articulan responsabilidad compartida y dominio conceptual individual.

\section{Marco conceptual del interaprendizaje en ambientes virtuales}

\subsection{Noción de interaprendizaje}

El \textbf{interaprendizaje} designa la construcción conjunta del conocimiento a
partir del intercambio de ideas, experiencias y perspectivas entre pares. No es un
método para repartir responsabilidades, sino una forma de aprender en la que la
interacción modifica lo que cada persona sabe. En entornos digitales, esa
construcción se sostiene cuando los recursos y la comunicación en línea se orientan a
fortalecer el aprendizaje colaborativo y autónomo, y no solo a transmitir contenidos
\citep{bertely2011interaprendizajes}. Así entendido, el interaprendizaje es el
resultado visible del \textbf{socioaprendizaje}: el aprendizaje social que ocurre
gracias a la interacción sostenida dentro de una comunidad.

\subsection{Principios y condiciones}

El interaprendizaje se reconoce por algunas condiciones que lo hacen posible: la
\textbf{escucha activa}, la disposición a transformar las ideas propias, la
corresponsabilidad en la construcción del conocimiento y la existencia de reglas de
participación que permitan una diversidad real de voces. Las metodologías activas para
el aprendizaje virtual muestran que estas condiciones no surgen de manera espontánea:
requieren estructurar la interacción, distribuir la palabra y sostener el seguimiento
del proceso, de modo que la colaboración se convierta en aprendizaje y no en simple
coordinación de entregas \citep{diazAlonso2011construccion}. En síntesis, tres
condiciones resultan decisivas:

\begin{itemize}[leftmargin=1.6em,itemsep=0.3em]
  \item \textbf{Construcción conjunta}, no reparto de tareas: las aportaciones se
  integran y se enriquecen mutuamente.
  \item \textbf{Transformación por interacción}, no acumulación: las ideas cambian al
  contrastarse con las de otras personas.
  \item \textbf{Diversidad incorporada}, no neutralizada: las diferencias de voz,
  lengua y perspectiva se leen como recurso y no como obstáculo.
\end{itemize}

\subsection{Prácticas dialógicas frente a prácticas monológicas}

La distinción central para participar bien es la que separa lo dialógico de lo
monológico. Una \textbf{práctica dialógica} abre el intercambio: pregunta, integra y
permite que el sentido se construya entre varias personas. Una \textbf{práctica
monológica}, en cambio, cierra el intercambio: impone una sola voz, simula acuerdo o
reduce la participación a un trámite. La comunicación intercultural advierte que en
toda interacción las personas no solo transmiten datos, sino que narran desde sus
marcos de autoridad, memoria y pertenencia; ignorarlo convierte la diferencia en
ruido y empobrece la construcción conjunta \citep{pech2016manual}. El interaprendizaje
se distorsiona, entonces, cuando la interacción se reduce a repartir tareas, repetir
opiniones, imponer una voz dominante o tratar la diversidad como un obstáculo en lugar
de un recurso.

\subsection{Contribución a la wiki colaborativa}

Con el marco anterior como sustento, se reproduce a continuación la contribución
publicada en la wiki de la semana, cuyo tema es aquello que el interaprendizaje no es
y cómo evitarlo. La contribución sigue la lógica del socioaprendizaje: aporta una
práctica que obstaculiza la construcción conjunta, explica por qué representa un
problema y ofrece una alternativa para transformarla, además de ampliar —sin
modificarla— una aportación previa de otra persona.

\begin{wikibox}[title={Contribución publicada en la wiki colaborativa}]
\hfill\wikiCopyButton{wiki-participacion-Actividad-3.txt}\\[-0.5em]
\phantomsection\label{wiki-participacion}
\wikimeta{Lo que el interaprendizaje NO es y cómo evitarlo}{Martin Jonathan de la Cruz}{\today}{Wiki colaborativa, aula virtual UnADM}

\textbf{Aportación: el monopolio de la palabra disfrazado de liderazgo.}\par
\medskip
Una práctica frecuente que contradice el interaprendizaje es que una sola persona
concentre la palabra y las decisiones del equipo, presentándolo como ``liderazgo'' o
``eficiencia''. En un ambiente virtual esto se reconoce cuando alguien redacta el
producto casi completo, fija las conclusiones y deja a las demás personas el papel de
aprobar lo ya hecho.

\textit{Por qué es un problema.} El monopolio de la palabra rompe la condición básica
del interaprendizaje: la construcción conjunta. Si las ideas no se contrastan ni se
transforman por la interacción, el producto reproduce una sola voz y deja de constituir
aprendizaje compartido \citep{bertely2011interaprendizajes}. A mi juicio, el problema no
se limita a una distribución desigual del trabajo: también revela que la interacción no
fue estructurada para producir aprendizaje, pues colaborar exige organizar la
participación y no solo coordinar entregas \citep{diazAlonso2011construccion}. Además,
invisibiliza la diversidad de perspectivas y lenguas del grupo, que debería enriquecer
el intercambio; tratarla como ruido empobrece la construcción conjunta \citep{pech2016manual}.

\textit{Alternativa para transformarla.} Una salida concreta es estructurar la
participación con roles rotativos y aportes mínimos verificables: cada persona
introduce una idea propia antes de comentar la ajena, y quien coordina sintetiza sin
sustituir. En el trabajo académico a distancia esto se traduce en fijar turnos de
edición en el documento compartido; en el campo jurídico, en un equipo que analiza un
caso, equivale a que cada integrante aporte primero su lectura del problema antes de
cerrar una estrategia común, para que la decisión final integre varias perspectivas y
no la de una sola voz.

\medskip
\textbf{Complemento a una aportación previa.} Ampliando —sin modificarla— una
aportación anterior que señalaba el ``consenso aparente'' como práctica dañina, añado
un ejemplo del ámbito jurídico: cuando un equipo aprueba rápido una conclusión ``para
avanzar'', suele ocultar desacuerdos que resurgen al final. Una forma de prevenirlo es
pedir, antes de cerrar, al menos una objeción o matiz por persona; así el acuerdo se
construye y no solo se simula.

\medskip
\textbf{Para seguir construyendo entre todas y todos:} ¿qué señales tempranas les
ayudan a distinguir, en una actividad en línea, entre una colaboración real y un
reparto de tareas que solo aparenta trabajo conjunto?
\end{wikibox}

Leída en conjunto, la contribución muestra que reconocer una práctica monológica no
basta: hace falta proponer un cambio en la dinámica. El monopolio de la palabra se
corrige distribuyendo la participación, y el consenso aparente, provocando el
disenso antes de cerrar. En ambos casos, la diversidad de voces deja de ser un
obstáculo para volverse la condición del aprendizaje compartido.

\subsection{Cuestionario: comprender, construir y distinguir}

El cuestionario sumativo aplica la técnica de estudio de caso y consta de diez
reactivos de opción múltiple. Cada respuesta correcta vale un punto; el intento es
único y dispone de treinta minutos. Los primeros reactivos verifican la noción, las
condiciones y las habilidades del interaprendizaje; los últimos trasladan ese marco a
un equipo que colabora a distancia y prepara material jurídico para una comunidad.
Esta progresión es congruente con la evidencia de que los entornos virtuales pueden
favorecer aprendizaje activo, autónomo y colaborativo cuando la interacción se
organiza con intención pedagógica \citep{cushpa2022pacie,silva2015metodologias}.

La tabla conserva las respuestas registradas en la nota de apoyo del 25 de julio de
2026. Dicha nota contiene los reactivos 1 a 8 y 10, pero omite el número 9; por
integridad académica, ese vacío se identifica expresamente y no se sustituye por una
pregunta inventada.

El patrón de respuestas permite distinguir con precisión la colaboración aparente del
interaprendizaje. Organizar tareas y utilizar plataformas son condiciones operativas,
pero solo hay construcción conjunta cuando se intercambian saberes, se revisa el
producto completo, se distribuyen las decisiones y se adapta la comunicación a las
personas destinatarias. En el caso jurídico, este último punto es decisivo: dialogar
con la comunidad mejora tanto la pertinencia del material como la legitimidad del
proceso mediante el cual se elabora.

\clearpage
\pagestyle{empty}
\begin{landscape}
\thispagestyle{empty}
\begingroup
\fontsize{7.25}{8.1}\selectfont
\setlength{\tabcolsep}{2.8pt}
\renewcommand{\arraystretch}{1.12}
\begin{longtable}{@{}>{\centering\arraybackslash}p{0.045\linewidth}>{\raggedright\arraybackslash}p{0.37\linewidth}>{\raggedright\arraybackslash}p{0.25\linewidth}>{\raggedright\arraybackslash}p{0.27\linewidth}@{}}
\caption{Cuestionario sobre interaprendizaje y colaboración académica a distancia}\label{tab:cuestionario-interaprendizaje-a3}\\
\toprule
\textbf{No.} & \textbf{Pregunta} & \textbf{Respuesta correcta} & \textbf{Justificación} \\
\midrule
\endfirsthead
\toprule
\textbf{No.} & \textbf{Pregunta} & \textbf{Respuesta correcta} & \textbf{Justificación} \\
\midrule
\endhead
\bottomrule
\endfoot
\bottomrule
\endlastfoot
% Opciones del reactivo 1 (comentadas, no visibles):
% a. Es una estrategia de organización que distribuye equitativamente una actividad entre varias personas.
% b. Es una forma de enseñanza en la que una persona con mayor experiencia transmite sus conocimientos al resto.
% c. Es una modalidad de estudio individual sin necesidad de interactuar con otras personas.
% d. Es un proceso de construcción colectiva del conocimiento, basado en la interacción recíproca y el diálogo.
1 & ¿Cuál de las siguientes opciones explica de manera más completa qué es el interaprendizaje? & \textbf{d.} Es un proceso de construcción colectiva del conocimiento, basado en la interacción recíproca y el diálogo. & El interaprendizaje se define por la construcción colectiva y el diálogo recíproco \citep{bertely2011interaprendizajes}. \\
\midrule
% Opciones del reactivo 2 (comentadas, no visibles):
% a. Permite tomar decisiones con rapidez aceptando sin objeciones la opinión de la mayoría.
% b. Favorece el contraste de interpretaciones, la construcción de argumentos y la búsqueda conjunta de alternativas.
% c. Permite reemplazar el estudio individual de las normas por explicaciones del grupo.
% d. Reduce la responsabilidad individual de cada estudiante.
2 & ¿Qué implicación puede tener el interaprendizaje en la formación académica de las y los estudiantes? & \textbf{b.} Favorece el contraste de interpretaciones, la construcción de argumentos y la búsqueda conjunta de alternativas. & El interaprendizaje potencia el diálogo argumentativo y la búsqueda colectiva. \\
\midrule
% Opciones del reactivo 3 (comentadas, no visibles):
% a. Consiste en traducir literalmente la información sin modificar su contenido.
% b. Favorece el reconocimiento de diferentes saberes, lenguas y experiencias, ampliando las formas de interpretar la realidad mediante el diálogo.
% c. Procura reducir las diferencias culturales para que todos usen las mismas formas.
% d. Permite que cada grupo conserve sus conocimientos por separado.
3 & ¿Cómo se relaciona el interaprendizaje con la diversidad cultural y lingüística? & \textbf{b.} Favorece el reconocimiento de diferentes saberes, lenguas y experiencias, ampliando las formas de interpretar la realidad mediante el diálogo. & El interaprendizaje reconoce y valora la diversidad mediante el diálogo \citep{pech2016manual}. \\
\midrule
% Opciones del reactivo 4 (comentadas, no visibles):
% a. Gestiona información de manera crítica, actúa con autonomía, escucha, argumenta, colabora y ajusta sus decisiones.
% b. Participa solo cuando recibe instrucciones y evita los desacuerdos.
% c. Domina los contenidos, defiende sus opiniones y trabaja sin depender de otros.
% d. Utiliza herramientas digitales y se concentra solo en su parte asignada.
4 & ¿Cuál de los siguientes perfiles corresponde mejor a una persona estudiante que aplica el interaprendizaje? & \textbf{a.} Gestiona información de manera crítica, actúa con autonomía, escucha, argumenta, colabora y ajusta sus decisiones. & El perfil colaborativo, crítico y abierto a nuevas perspectivas es propio del interaprendizaje. \\
\midrule
% Opciones del reactivo 5 (comentadas, no visibles):
% a. La libertad de trabajar sin acuerdos previos ni seguimiento.
% b. Un ambiente de confianza y respeto, relaciones horizontales, participación comprometida y espacios organizados para la interacción.
% c. Comunicación frecuente, misma cantidad de trabajo y obligación de acuerdo rápido.
% d. Una plataforma tecnológica, una persona que controle las decisiones y distribución individual de tareas.
5 & ¿Qué conjunto de condiciones favorece el desarrollo del interaprendizaje? & \textbf{b.} Un ambiente de confianza y respeto, relaciones horizontales, participación comprometida y espacios organizados para la interacción. & La confianza, horizontalidad y participación son condiciones clave del interaprendizaje. \\
\midrule
% Opciones del reactivo 6 (comentadas, no visibles):
% a. Porque debe realizarse únicamente por videollamadas.
% b. Porque la distribución de tareas impide la responsabilidad individual.
% c. Porque cada integrante trabaja en una sección aislada, sin intercambio de saberes ni revisión conjunta.
% d. Porque todos deberían redactar al mismo tiempo cada apartado.
6 & ¿Por qué la organización inicial del equipo todavía no representa una experiencia completa de interaprendizaje? & \textbf{c.} Porque cada integrante trabaja en una sección aislada del documento, sin intercambio de saberes ni revisión conjunta del producto. & Faltan la interacción y revisión conjunta propias del interaprendizaje. \\
\midrule
% Opciones del reactivo 7 (comentadas, no visibles):
% a. Acordar una ruta de trabajo, combinar espacios sincrónicos y asincrónicos que integren los aportes de quienes no logran conectarse, y revisar el documento de forma colectiva antes de entregarlo.
% b. Establecer una sola videollamada obligatoria sin alternativas y descartar aportaciones de quienes no se conecten.
% c. Permitir que la persona coordinadora tome todas las decisiones por su cuenta.
% d. Dejar que cada integrante elija su herramienta y entregue su parte sin coordinación.
7 & ¿Qué acción habría permitido al equipo sostener una colaboración a distancia más congruente con los principios del interaprendizaje? & \textbf{a.} Acordar una ruta de trabajo, combinar espacios sincrónicos y asincrónicos que integren los aportes de quienes no logran conectarse, y revisar el documento de forma colectiva antes de entregarlo. & Integrar aportes de todos y revisar colectivamente refleja el interaprendizaje. \\
\midrule
% Opciones del reactivo 8 (comentadas, no visibles):
% a. Agregar un glosario bilingüe elaborado solo por el equipo jurídico, sin consultar a la comunidad.
% b. Dialogar con personas de la comunidad, recuperar sus conocimientos y adaptar el lenguaje, los ejemplos y los formatos de comunicación.
% c. Conservar el lenguaje técnico y traducir literalmente solo algunos términos.
% d. Evitar cualquier referencia cultural para mantener el contenido neutral.
8 & ¿Qué decisión sería más congruente con el interaprendizaje al elaborar el material jurídico para la comunidad? & \textbf{b.} Dialogar con personas de la comunidad, recuperar sus conocimientos y adaptar el lenguaje, los ejemplos y los formatos de comunicación. & Consultar y adaptar según la comunidad refleja el diálogo intercultural del interaprendizaje. \\
\midrule
9 & Reactivo no conservado en la nota de apoyo. & \textbf{No verificable con la evidencia disponible.} & Se mantiene la trazabilidad del instrumento: atribuir una pregunta o respuesta no registrada vulneraría la integridad académica del producto. \\
\midrule
% Opciones del reactivo 10 (comentadas, no visibles):
% a. Programar más videollamadas sin modificar la distribución del trabajo ni los mecanismos de decisión.
% b. Recuperar las aportaciones de todas y todos, revisar y respetar los acuerdos, combinar canales de participación, integrar los conocimientos comunitarios y evaluar juntos el producto.
% c. Mantener la distribución inicial y pedir a la coordinadora que corrija las contradicciones.
% d. Encargar la versión final a quien tenga mejor redacción y traducir después los términos.
10 & ¿Qué plan permitiría al equipo reorganizarse conforme a los principios del interaprendizaje? & \textbf{b.} Recuperar las aportaciones de todas y todos los participantes, revisar y respetar los acuerdos, combinar canales de participación, integrar los conocimientos comunitarios y evaluar juntos el producto. & Integrar aportes, respetar acuerdos y evaluar en conjunto es coherente con el interaprendizaje. \\
\end{longtable}
\endgroup
\end{landscape}
\clearpage
\pagestyle{fancy}

\clearpage
\section{Conclusión}

A mi juicio, comprender el interaprendizaje es, sobre todo, aprender a distinguir la
construcción conjunta real de la colaboración aparente. Las prácticas que lo
obstaculizan —el monopolio de la palabra, el consenso aparente, la participación
decorativa o la invisibilización de las diferencias— comparten un mismo rasgo:
sustituyen el diálogo por una sola voz y convierten la diversidad en silencio.

Por ello sostengo que el valor de los dos productos está en su complementariedad. La
wiki exige aportar de modo que el conocimiento se construya entre varias personas; el
cuestionario exige distinguir individualmente, ante un caso, qué decisiones sostienen
esa construcción. Ese criterio —trabajar por distinciones y no por definiciones
memorizadas— vuelve útil el concepto tanto en el trabajo académico a distancia como
en el ejercicio del Derecho, donde decidir en conjunto, justificar alternativas y
escuchar la diferencia son condiciones de una mejor
justicia.\footnote{Para redactar y organizar este documento se utilizó una
herramienta de inteligencia artificial con el propósito de ordenar ideas y revisar la
redacción; su uso no sustituyó la comprensión de las lecturas, la selección y
aplicación de las fuentes ni la elaboración de los productos, que corresponden al
autor.}

\clearpage
\nocite{bertely2011interaprendizajes,diazAlonso2011construccion,pech2016manual,cushpa2022pacie,silva2015metodologias}
\bibliography{interaprendizaje-en-ambientes-virtuales}

\end{document}
